\section{Self-Organising Map}
\label{sec:SOM}
A \ac{SOM} is an unsupervised machine learning tecnique used to produce low-dimensionality representation of a higher-dimensionality dataset preserving topological structures of data that distinguish different types of said data. Specifically, a \ac{SOM} is a type of artificial neural network trained using competitive learning rather than error-correction learning. The goal in these auto-organised learning maps is to train different parts of the \ac{SOM} to reply to different input patterns: this specific design has been formulated in the 1980s by Teuvo Kohonen on the basis of how the human brain manages different information in different parts of the cerebral cortex.

Starting from weights vectors initialised to casual or very small numbers, the algorithm exploits competitive learning: when a new training sample is received by the network, a distance metric is computed from all the weights vectors. The most alike weights vector, named \ac{BMU} and all the neurons near to ir in the \ac{SOM} net are brought closer with one another and weights are updated.